\section{Conclusiones y Trabajo Futuro}
\label{sec:conclusiones}

Este proyecto valid\'o la integraci\'on completa entre el stack cu\'antico-computacional \textit{Qiskit} y modelos m\'etricos cu\'anticos operados f\'isicamente en un procesador IBM Heron r2 (\texttt{ibm\_fez}, 156 qubits). A modo de s\'intesis:

\begin{itemize}
    \item \textbf{Verificaci\'on del hardware:} La preparaci\'on del estado de Bell alcanz\'o una fidelidad de 0.9834, confirmando la alta calidad del hardware para operaciones de entrelazamiento b\'asico.
    \item \textbf{Impacto dual cuantificado:} Se demostr\'o que la degradaci\'on del rendimiento de la QSVC en hardware real responde a dos factores simult\'aneos y acumulativos: (i) la reducci\'on muestral de 105 a 13 muestras de entrenamiento provoc\'o una ca\'ida de 77.78\% a 33.33\% en simulador ideal ($-$44 pp), y (ii) el ruido NISQ redujo adicionalmente el resultado de 33.33\% a 11.11\% en hardware ($-$22 pp).
    \item \textbf{Verificaci\'on Variacional:} El VQC requiere, como m\'inimo, 300 iteraciones de COBYLA para escapar de m\'inimos locales tempranos, alcanzando un m\'aximo de 60\% de accuracy ---a\'un 33 pp por debajo de la frontera cl\'asica---.
    \item \textbf{Ineficacia de TREX:} La mitigaci\'on de errores de lectura (TREX) no logr\'o recuperar rendimiento bajo la degradaci\'on combinada de escasez muestral y ruido de compuerta.
    \item \textbf{Anomal\'ia MNIST:} La clasificaci\'on binaria en hardware (88.89\%) super\'o al simulador (77.78\%), fen\'omeno atribuible a la baja resoluci\'on estad\'istica ($\pm$11.1\% por muestra) y no a una ventaja cu\'antica genuina.
\end{itemize}

\subsection*{Proyecciones a Futuro}
La investigaci\'on requiere tres l\'ineas de avance complementarias:
\begin{enumerate}
    \item \textbf{Escala muestral:} Migrar a planes de acceso premium (IBM Premium Access Programs) que permitan inyectar la totalidad de los datos en el kernel f\'isico, eliminando el factor de escasez muestral y aislando el efecto puro del ruido.
    \item \textbf{Mitigaci\'on avanzada:} Implementar t\'ecnicas que operen sobre el ruido de compuerta, como \textbf{ZNE} (Zero Noise Extrapolation) o \textbf{PEC} (Probabilistic Error Cancellation), junto a la primitiva \textit{EstimatorV2}, para evaluar su capacidad de recuperaci\'on sobre kernels cu\'anticos \citep{larocca2025}.
    \item \textbf{Benchmark justo:} Replicar la evaluaci\'on sistem\'atica de \citet{bowles2024} con las 160 bases de datos, utilizando hardware real con datasets completos, para determinar bajo qu\'e condiciones espec\'ificas los modelos cu\'anticos pueden ser competitivos frente a los cl\'asicos \citep{huang2021, biamonte2017}.
\end{enumerate}
